% slide-ready summary for the lecture box 'Research from our group' (max 6 lines); generated by the notebook
\begin{itemize}\setlength{\itemsep}{2pt}
  \item \textbf{Avertizarea explicabilă privind vîrfurile de preț} (Andrei, Pele, Cramer, Amza și Diaconu, 2026, PICBE): prețul anterior și consumul rezidual determină vîrfurile (SHAP, Tabelul~5); sistemul are o sensibilitate de 0,722 și un scor F1 de 0,569 (Tabelul~7).
  \item Replicarea noastră (2023--oct.\ 2026, numai cu informația de la închiderea porții): factorii principali sînt prețul, aceeași oră ieri și prețul, aceeași oră săptămîna trecută; XGBoost obține un AUC de 0,95 (cu previziune perfectă a variabilelor fundamentale: 0,89).
  \item Totuși, vîrfurile au devenit zilnice (22\,\% din orele de test, față de 10\,\% în antrenare): regula „vîrf ieri” atinge F1 0,77, față de 0,63 pentru XGBoost --- recalibrați pragurile.
  \item Decizia: reacționați atunci cînd probabilitatea unui vîrf depășește $C/L$, raport stabilit din propriile costuri. Vedeți și modelele fundamentale pentru riscul extrem (Jheng et al., 2026, ECECSR).
\end{itemize}
